
\documentclass{article}
\usepackage{colortbl}
\usepackage{makecell}
\usepackage{multirow}
\usepackage{supertabular}

\begin{document}

\newcounter{utterance}

\centering \large Interaction Transcript for game `codenames', experiment `frequency\_low', episode 3 with reflexion{-}Qwen3.5{-}27B.
\vspace{24pt}

{ \footnotesize  \setcounter{utterance}{1}
\setlength{\tabcolsep}{0pt}
\begin{supertabular}{c@{$\;$}|p{.15\linewidth}@{}p{.15\linewidth}p{.15\linewidth}p{.15\linewidth}p{.15\linewidth}p{.15\linewidth}}
    \# & $\;$A & \multicolumn{4}{c}{Game Master} & $\;\:$B\\
    \hline

    \theutterance \stepcounter{utterance}  
    & & & \multicolumn{2}{p{0.3\linewidth}}{
        \cellcolor[rgb]{0.9,0.9,0.9}{
            \makecell[{{p{\linewidth}}}]{
                \texttt{\tiny{[GM$|$GM]}}
                {'hidden': {'team': ['unchallenging', 'hypersensitivity'], 'innocent': ['bandmaster', 'ladylove', 'customizer', 'defenseman', 'narrowness', 'hotelkeeper', 'neckband'], 'opponent': [], 'assassin': ['seacoast']}, 'revealed': {'team': {'team': [], 'innocent': [], 'opponent': [], 'assassin': []}, 'opponent': {'team': [], 'innocent': [], 'opponent': [], 'assassin': []}}}
            }
        }
    }
    & & \\ \\

    \theutterance \stepcounter{utterance}  
    & & \multicolumn{4}{p{0.6\linewidth}}{
        \cellcolor[rgb]{0.9,0.9,0.9}{
            \makecell[{{p{\linewidth}}}]{
                \texttt{\tiny{[P1$\langle$GM]}}
                \texttt{You have attempted this same task before and did not succeed. These are the plans you wrote after those attempts.} \\
\texttt{Trial \#1: New plan: I failed because I gave clues for only one word at a time in the final turns, which was inefficient and allowed the opponent to catch up. Specifically, I missed the opportunity to group "unchallenging" and "hypersensitivity" together earlier. Although "SOFT" was a risky clue that led to an opponent guess, I should have used a safer, more specific clue like "SENSITIVE" or "REACTIVE" to target "hypersensitivity" while potentially linking to "unchallenging" (as in "not sensitive to challenge") or simply accepting the risk to clear two words at once. In the future, I will prioritize grouping remaining words into multi{-}word clues to maximize points per turn, even if the association is slightly more abstract, rather than playing it safe with single{-}word clues that waste turns. I will also avoid clues that are too broad (like "SOFT") if they strongly overlap with opponent words like "liquidness" or "emotionality".} \\
\\ 
\texttt{Trial \#2: New plan: I failed because I only targeted "hypersensitivity" with "REACTIVE," leaving "unchallenging" unguessed and wasting a turn. I missed the semantic link that "unchallenging" implies a lack of difficulty or a "soft" task, while "hypersensitivity" implies an over{-}reaction to stimuli. A better strategy is to find a single word that bridges both concepts, such as "EASY" (relating to "unchallenging" and potentially "easy to trigger" for hypersensitivity) or "SENSITIVE" (directly linking to hypersensitivity and perhaps "sensitive" to challenge). Specifically, I should have attempted a clue like "SENSITIVE" or "EASY" with TARGETS: unchallenging, hypersensitivity to clear both words in one turn, rather than splitting them. If a perfect bridge is too risky, I must prioritize the word with the strongest, most unique association first, but always aim for multi{-}word clues to maximize efficiency.} \\
\\ 
\\ 
\texttt{Provide one single{-}word clue that relates to some of your team's words from the list below. You can choose to how many and to which words your clue relates to.} \\
\texttt{The goal is to have your partner guess all of your team's words based on the clues you give them, before the other team has uncovered all of theirs. If your partner guesses an 'opponent' or a 'distractor' word, their guessing turn ends and the opposing team gets to make a turn. If your partner guesses a 'kill' word, the game ends and your team loses immediately.} \\
\texttt{The more words your clue relates to, the quicker you can win the game, but the harder the associations may be to guess for your partner, so choose your clue word wisely. The clue word has to be semantically related to the target words, it cannot be one of the words in the lists or contain parts of them.} \\
\\ 
\texttt{Always give your single{-}word clue and your comma{-}separated list of related target words in the following format and make your answers as short as possible, never include any other text than is required in this form:} \\
\\ 
\texttt{CLUE: <WORD>} \\
\texttt{TARGETS: <TARGETS>} \\
\\ 
\texttt{Your team words are: unchallenging, hypersensitivity.} \\
\texttt{Your opponent's team words are: .} \\
\texttt{Distractor words are: bandmaster, ladylove, customizer, defenseman, narrowness, hotelkeeper, neckband.} \\
\texttt{The kill words are: seacoast.} \\
            }
        }
    }
    & \\ \\

    \theutterance \stepcounter{utterance}  
    & \multicolumn{4}{p{0.6\linewidth}}{
        \cellcolor[rgb]{0.9,0.9,0.9}{
            \makecell[{{p{\linewidth}}}]{
                \texttt{\tiny{[P1$\rangle$GM]}}
                \texttt{CLUE: EASY} \\
\texttt{TARGETS: unchallenging, hypersensitivity} \\
            }
        }
    }
    & & \\ \\

    \theutterance \stepcounter{utterance}  
    & & & \multicolumn{2}{p{0.3\linewidth}}{
        \cellcolor[rgb]{0.9,0.9,0.9}{
            \makecell[{{p{\linewidth}}}]{
                \texttt{\tiny{[GM$|$GM]}}
                \texttt{easy} \\
            }
        }
    }
    & & \\ \\

    \theutterance \stepcounter{utterance}  
    & & & \multicolumn{2}{p{0.3\linewidth}}{
        \cellcolor[rgb]{0.9,0.9,0.9}{
            \makecell[{{p{\linewidth}}}]{
                \texttt{\tiny{[GM$|$GM]}}
                ['unchallenging', 'hypersensitivity']
            }
        }
    }
    & & \\ \\

    \theutterance \stepcounter{utterance}  
    & & & \multicolumn{2}{p{0.3\linewidth}}{
        \cellcolor[rgb]{0.9,0.9,0.9}{
            \makecell[{{p{\linewidth}}}]{
                \texttt{\tiny{[GM$|$GM]}}
                {'word': 'unchallenging', 'assignment': 'team'}
            }
        }
    }
    & & \\ \\

    \theutterance \stepcounter{utterance}  
    & & & \multicolumn{2}{p{0.3\linewidth}}{
        \cellcolor[rgb]{0.9,0.9,0.9}{
            \makecell[{{p{\linewidth}}}]{
                \texttt{\tiny{[GM$|$GM]}}
                {'word': 'hypersensitivity', 'assignment': 'team'}
            }
        }
    }
    & & \\ \\

    \theutterance \stepcounter{utterance}  
    & & & \multicolumn{2}{p{0.3\linewidth}}{
        \cellcolor[rgb]{0.9,0.9,0.9}{
            \makecell[{{p{\linewidth}}}]{
                \texttt{\tiny{[GM$|$GM]}}
                \texttt{opponent has won} \\
            }
        }
    }
    & & \\ \\

    \theutterance \stepcounter{utterance}  
    & \multicolumn{4}{p{0.6\linewidth}}{
        \cellcolor[rgb]{0.95,0.95,0.95}{
            \makecell[{{p{\linewidth}}}]{
                \texttt{\tiny{[P2$\rightarrow$P1]}}
                \texttt{Post{-}episode Reflexion (not sent to GM):} \\
\texttt{New plan: I failed because the clue "EASY" was semantically valid for "unchallenging" but too weak and ambiguous for "hypersensitivity," likely causing my partner to guess a distractor or fail to connect the two. "Hypersensitivity" implies an *exaggerated* or *intense* reaction, which is the opposite of "easy." I should have chosen a clue that captures the core concept of "sensitivity" or "reaction" which applies to both: "unchallenging" can be interpreted as a situation requiring no "sensitivity" (or a "low" sensitivity threshold), while "hypersensitivity" is an extreme form of it. A better strategy is to use a clue like "SENSITIVE" or "REACTION" and explicitly list both words as targets, trusting the partner to make the logical leap that "unchallenging" relates to the *absence* or *low level* of the clue's concept, or to accept that "SENSITIVE" is the primary link and hope the partner guesses "unchallenging" as the antonym or related state. Specifically, I will use CLUE: SENSITIVE with TARGETS: unchallenging, hypersensitivity, as "sensitive" is the direct root of the second word and the conceptual opposite of the first, creating a strong, albeit slightly abstract, bridge. If that feels too risky, I will prioritize the most distinct word first with a precise clue like "OVERREACT" for "hypersensitivity" to secure one point safely before attempting a bridge.} \\
            }
        }
    }
    & & \\ \\

\end{supertabular}
}

\end{document}
